We study a finite, self-contained question inside \(\Sfour\): when a subset of the 24 one-line words is viewed either as a set of linear extensions of a strict poset on \([4]\), or as a subgroup row of \(\Sfour\), the two structures need not coincide. The paper first fixes the indexing dictionary used for the finite layer. One-line words, local \(A_3\) prefix flags, and \(4\times4\) permutation-matrix supports are canonically bijective, while prefix-support vectors are exactly the cardinality vectors of unordered set-prefix shadows.

For every nonempty \(X\subseteq\Sfour\), the common-precedence relation records the ordered pairs of labels whose relative order is forced in every word of \(X\). We prove that its transitive closure is a strict poset \(P_X\), and that \(X\) is an exact poset cone precisely when \(X=\Lin(P_X)\). The selected finite data package records bounded package counts and replay checks accompanying this comparison. The mathematical subset count used below is 219 distinct poset-cone subsets from strict labeled posets on \([4]\); the other selected counts are artifact-row accounting for the packaged rows.

The main witness separation is explicit. A five-word N-poset row is an exact poset cone but not a subgroup row, while the locked \(\Dfour\) and \(\Vfour\) rows are subgroup rows whose common-precedence closure is all of \(\Sfour\), so they are not exact poset cones. Thus the displayed rows witness that exact poset-cone status is not equivalent to subgroup-row status. Left and right coset counts for the subgroup rows are recorded as induced finite group bookkeeping, not as a separate structural equivalence theorem. The paper does not claim a classification of all subsets of \(\Sfour\), an all-tiler theorem, a geometric realization theorem, or a general \(\Sn\) theorem.